\documentclass[11pt]{article}
\usepackage{ochamath}
\subjectcolor{2F5DA8}
\setsheet{線形代数・電気系院試補充}{行列式の計算法とブロック行列　演習}{https://university-math-crj.pages.dev/linear-algebra/determinant-methods/}
\namelinetrue
\begin{document}
\makesheethead{問題}
自作問題。本文の例題とは数値や設定が異なります。条件・途中式・検算を示してください。
\problem[消去と行列式]
$A=\begin{pmatrix}3&1&0\\1&2&1\\0&1&2\end{pmatrix}$ の行列式を、行列式を変えない行操作で求め、余因子展開で検算せよ。
\workspace{45mm}
\problem[随伴行列]
$A=\begin{pmatrix}3&1\\2&1\end{pmatrix}$ の随伴行列と逆行列を求め、左右の積を確かめよ。
\workspace{45mm}
\newpage
\problem[クラメルの公式]
$3x+y=5,\ 2x+y=4$ をクラメルの公式で解き、代入して検算せよ。
\workspace{45mm}
\problem[Schur補行列]
$M=\begin{pmatrix}3&0&1\\0&2&2\\1&2&4\end{pmatrix}$ の行列式を左上の2次ブロックで求めよ。
\workspace{45mm}
\newpage
\problem[内部変数の消去]
$3x_1+y=4,\ 2x_2+2y=6,\ x_1+2x_2+4y=9$ をSchur補行列で解け。
\workspace{45mm}
\problem[行列式と相似]
正則な $P$ と $n$ 次行列 $A$ について $\det(P^{-1}AP)=\det A$ と $\det(cA)=c^n\det A$ を証明せよ。
\workspace{45mm}
\end{document}
